% Capítulo 3: el canal analógico (fuente IEPE, sumador, DAC y lazo de DC)
\chapter[El canal analógico]{El canal analógico: fuente IEPE, sumador y lazo de DC}\label{cap:analog}

Este es el corazón del sonómetro. Todo lo demás de la placa (fuentes conmutadas, ADC, NUCLEO) existe para que este
canal funcione bien. Aquí el micrófono se alimenta, su señal se separa de la tensión de polarización y queda lista
para que el ADC la digitalice.

En este capítulo vas a ver:
\begin{itemize}
  \item qué tiene que hacer el canal y cómo está dividido en bloques;
  \item cómo funciona la fuente de corriente IEPE con el LT3092 y de dónde salen los \SI{2.81}{\milli\ampere} y
        \SI{1.51}{\milli\ampere} de cada canal (y el escalón de los puentes JP1/JP2);
  \item el sumador inversor que resta la polarización del micrófono: la ecuación deducida paso a paso y por qué la
        ganancia es \num{0.1};
  \item cómo el DAC MCP4725, su filtro y el seguidor generan la referencia \net{VSHIFT}, y cómo el lazo de DC del
        firmware la ajusta;
  \item por qué el sumador necesitaba un condensador de realimentación (CF = \SI{10}{\pico\farad}) para no oscilar;
  \item qué cambió al pasar del TLV9062 al TLV9152 y las dos comprobaciones que hubo que hacer;
  \item cómo el canal 2 copia al canal 1, y cómo se alimenta y desacopla todo esto.
\end{itemize}

Un detalle importante antes de empezar: \textbf{la topología de este canal no la inventamos durante el ruteo}. La
fuente IEPE, el sumador, el DAC y el seguidor del canal 1 vienen tal cual del proyecto Altium del equipo
(\archivo{Sonometro_2CH}); el canal 2 es una copia pieza por pieza (\archivo{tools/design.py}, tablas
\texttt{MIRROR} y \texttt{NETMAP}). Lo que se agregó o cambió en la Rev B (el condensador CF, el TLV9152, los filtros
del lazo, el escalón de corriente) se dice explícitamente en cada sección. Las conexiones de todas las figuras se
comprobaron contra \archivo{tools/target.json}, que es la lista de piezas y redes de la placa.

% =====================================================================================================
\section{Qué tiene que hacer el canal}\label{sec:c03-objetivo}

Un micrófono IEPE (\emph{Integrated Electronics Piezo-Electric}) tiene un preamplificador adentro y usa un solo cable
coaxial para dos cosas a la vez: por ese cable recibe una \textbf{corriente constante} que lo alimenta y por ese mismo
cable devuelve la \textbf{señal}. La tensión que ves en el centro del BNC es
\begin{equation}
  v_{\mathrm{IEPE}}(t) = V_B + S\cdot p(t),
  \label{eq:c03-iepe}
\end{equation}
donde $V_B$ es la tensión de polarización (la fija el propio micrófono, típicamente entre \SIrange{8}{14}{\volt}),
$S$ es la sensibilidad (por ejemplo \SI{50}{\milli\volt\per\pascal} en el RG-50 \doc{ROGA RG-50}) y $p(t)$ es la
presión sonora en pascales.

\begin{idea}[el problema en una frase]
La señal útil ($S\cdot p$) son unos milivoltios montados sobre unos \SI{12}{\volt} de DC. El ADC (PCM1808) solo acepta
\SI{3}{\volt} pico a pico centrados en su propia referencia \doc{SLES177B p.6}. Hay que quitar los \SI{12}{\volt},
achicar la señal para que quepa el peor caso y entregarla a un nivel de DC conocido.
\end{idea}

La figura~\ref{fig:c03-bloques} muestra cómo se reparte ese trabajo. Cada bloque tiene su sección en este capítulo.

\begin{figure}[htbp]
\centering
\begin{tikzpicture}[
  blq/.style={draw,rounded corners=2pt,align=center,font=\small,minimum height=11mm,inner sep=4pt,fill=white},
  sen/.style={-{Stealth[length=2.2mm]},thick},
  lbl/.style={font=\scriptsize\ttfamily,fill=white,inner sep=1pt}]
  \node[blq,fill=cobre!8] (fuente) at (0,0) {Fuente de corriente\\LT3092 (U1)};
  \node[blq,fill=azulidea!6] (mic) at (0,-2.4) {Micrófono IEPE\\(J2, BNC)};
  \node[blq,fill=verdeplaca!8] (sum) at (5.2,-1.2) {Sumador inversor\\U3A\\$\scriptstyle 1{,}1\,V_{\mathrm{SHIFT}}-0{,}1\,v_{\mathrm{IEPE}}$};
  \node[blq] (adc) at (10.2,-1.2) {C9 $\to$ PCM1808\\(ADC, cap.~\ref{cap:digital})};
  \node[blq] (mon) at (10.2,-4.0) {Divisor $\div 2$\\R15/R17, C8};
  \node[blq] (nuc) at (10.2,-6.4) {NUCLEO\\ADC PA0 + firmware};
  \node[blq] (dac) at (5.2,-6.4) {DAC MCP4725\\U2, I\textsuperscript{2}C 0x60};
  \node[blq] (rc) at (1.4,-6.4) {Filtro RC\\R12/C4, \SI{1}{\milli\second}};
  \node[blq] (seg) at (1.4,-4.2) {Seguidor\\U3B};
  \draw[sen] (fuente) -- node[lbl,left] {IEPE} (mic);
  \draw[sen] (0,-1.2) -- node[lbl,above] {IEPE (R13…R18, \SI{1}{\mega\ohm})} (sum.west);
  \draw[sen] (sum) -- node[lbl,above] {REST\_OUT} (adc);
  \draw[sen] (8.0,-1.2) |- (mon.west);
  \draw[sen] (mon) -- node[lbl,right] {DC\_MON} (nuc);
  \draw[sen] (nuc) -- node[lbl,above] {I2C} (dac);
  \draw[sen] (dac) -- node[lbl,above] {DAC\_RAW} (rc);
  \draw[sen] (rc) -- node[lbl,right] {DAC\_F} (seg);
  \draw[sen] (seg.north) |- node[lbl,pos=0.78,below] {VSHIFT} ([yshift=-3mm]sum.west);
  \node[font=\scriptsize,color=grisnota,align=left] at (5.4,-3.9) {lazo de DC:\\el firmware mueve VSHIFT\\hasta que REST\_OUT\\queda en su objetivo};
\end{tikzpicture}
\caption{Bloques del canal analógico del canal 1 con los nombres reales de las redes. El canal 2 es idéntico (U8,
U9, U10, con sufijo \texttt{\_SV} o \texttt{2} en las redes).}
\label{fig:c03-bloques}
\end{figure}

En palabras:
\begin{enumerate}
  \item La \textbf{fuente de corriente} empuja una corriente fija hacia el nodo \net{IEPE}. El micrófono la consume y
        fija ahí su tensión $V_B + S\,p(t)$.
  \item El \textbf{sumador inversor} U3A toma \net{IEPE} a través de \SI{1}{\mega\ohm} y le resta la referencia
        \net{VSHIFT}. Su salida \net{REST_OUT} es la señal atenuada diez veces, sobre un nivel de DC bajo y conocido.
  \item \net{REST_OUT} va al ADC por el condensador C9. Además, un divisor entre 2 la lleva al ADC de la NUCLEO como
        \net{DC_MON}, para que el firmware sepa cuál es el nivel de DC.
  \item El firmware calcula cuánto hay que mover \net{VSHIFT} y se lo manda al DAC por I\textsuperscript{2}C. El DAC,
        filtrado y copiado por el seguidor U3B, cierra el lazo.
\end{enumerate}

% =====================================================================================================
\section{La fuente de corriente IEPE con el LT3092}\label{sec:c03-fuente}

\subsection{Por qué el micrófono necesita una fuente de corriente}

\begin{idea}[fuente de voltaje frente a fuente de corriente]
Una fuente de voltaje mantiene fija la \emph{tensión} y deja que la corriente sea la que pida la carga. Una fuente de
corriente hace lo contrario: mantiene fija la \emph{corriente} y deja que la tensión sea la que imponga la carga. En
IEPE la tensión del cable \emph{es} la señal, así que la alimentación no puede imponerla: tiene que ser una fuente de
corriente, y el micrófono decide la tensión. Si pusieras una fuente de voltaje directo al micrófono, fijarías el nodo
y no habría señal (y probablemente lo dañarías: la regla de PCB es no aplicar nunca tensión sin limitar la corriente).
\end{idea}

Una fuente de corriente ideal tiene impedancia de salida infinita. Una real tiene una impedancia alta pero finita, y
eso importa: la señal del micrófono ve en paralelo esa impedancia, así que cuanto más alta, menos carga a la señal.
La fuente tampoco tiene que ser muy exacta en su valor (\SI{2.8}{\milli\ampere} o \SI{2.9}{\milli\ampere} da igual):
lo que importa es que sea \emph{estable}, que tenga poco ruido y que le alcance la tensión (\emph{compliance}). Esto lo
dice TI en su diseño de referencia TIDA-010249 \copia[AGENTS.md §4].

\subsection{Cómo funciona el LT3092}

El LT3092 (U1 en el canal 1, U8 en el canal 2) es una fuente de corriente programable con dos resistencias. Por dentro
tiene dos cosas \copia[ADI LT3092 Rev D]:
\begin{itemize}
  \item una fuente de corriente de precisión de \SI{10}{\micro\ampere} que sale por el pin \textsf{SET};
  \item un amplificador que obliga a que el pin \textsf{OUT} quede a la misma tensión que el pin \textsf{SET}.
\end{itemize}

Mira la figura~\ref{fig:c03-fuente}. Los \SI{10}{\micro\ampere} del pin SET pasan por $R_{\mathrm{SET}}$ y llegan al
nodo \net{IEPE}. Eso crea una caída $\SI{10}{\micro\ampere}\cdot R_{\mathrm{SET}}$ entre \net{SET_I} y \net{IEPE}.
Como el amplificador copia esa tensión en OUT, sobre $R_{\mathrm{OUT}}$ (entre \net{OUT_I} y \net{IEPE}) aparece
\emph{la misma caída}. Entonces por $R_{\mathrm{OUT}}$ circula
\begin{equation*}
  I_{R_{\mathrm{OUT}}} = \frac{\SI{10}{\micro\ampere}\cdot R_{\mathrm{SET}}}{R_{\mathrm{OUT}}}.
\end{equation*}
Las dos corrientes, la de SET y la de OUT, terminan en el nodo \net{IEPE}, así que la corriente total que recibe el
micrófono (más el poquito que se lleva el AFE, sección~\ref{sec:c03-carga}) es
\begin{equation}
  I_{\mathrm{IEPE}} = \SI{10}{\micro\ampere}\left(1 + \frac{R_{\mathrm{SET}}}{R_{\mathrm{OUT}}}\right).
  \label{eq:c03-lt3092}
\end{equation}

\begin{ojo}[el «+1» de la fórmula]
Es común ver la fórmula como $I = \SI{10}{\micro\ampere}\cdot R_{\mathrm{SET}}/R_{\mathrm{OUT}}$, sin el 1. La
diferencia son los \SI{10}{\micro\ampere} del propio pin SET, que también llegan a la carga. Con
$R_{\mathrm{SET}}/R_{\mathrm{OUT}} = 280$ el error de olvidarlo es de \SI{0.4}{\percent}: casi nada, pero las cuentas
de este capítulo lo incluyen para que coincidan con las notas del esquemático.
\end{ojo}

\begin{figure}[htbp]
\centering
\begin{circuitikz}[scale=0.74,transform shape,font=\small]
  % entrada y filtro
  \draw (0,4) node[left]{\net{24V_PROT}} to[nos,l=J1,o-] (2,4)
        to[R,l=R1,a=\SI{10}{\ohm}] (4.2,4) coordinate (sup);
  \draw (sup) node[above right]{\scriptsize\net{IEPE_SUP}};
  \draw (2,4) node[below,font=\scriptsize]{\net{IEPE_FEED}};
  \draw (4.2,4) to[short,*-] (4.2,3.6) to[C,l_={C2 \SI{1}{\micro\farad}}] (4.2,1.6) node[ground]{};
  \draw (4.2,4) -- (6.2,4) to[short,*-] (6.2,3.6) to[C,l_={C3 \SI{100}{\nano\farad}}] (6.2,1.6) node[ground]{};
  % LT3092
  \draw[thick] (8.0,2.6) rectangle (10.0,5.0);
  \node[align=center,font=\small] at (9.0,3.8) {U1\\LT3092};
  \draw (6.2,4) -- (8.0,4) node[right,font=\scriptsize]{IN};
  \draw (10.0,4.6) node[left,font=\scriptsize]{OUT} -- (11.0,4.6) coordinate (out);
  \draw (10.0,3.0) node[left,font=\scriptsize]{SET} -- (11.0,3.0) coordinate (set);
  \draw (out) node[above left,font=\scriptsize]{\net{OUT_I}};
  \draw (set) node[below left,font=\scriptsize]{\net{SET_I}};
  % ROUT = R2 || R3 y escalón
  \draw (out) to[R,l=R2,a=\SI{150}{\ohm},*-] (14.4,4.6) coordinate (iepe);
  \draw (out) -- (11.0,6.0) to[R,l=R3,a=\SI{150}{\ohm},*-] (14.4,6.0) -- (iepe);
  \draw (11.0,6.0) -- (11.0,7.4) to[R,l=R47,a=\SI{180}{\ohm}] (12.7,7.4) to[nos,l=JP1] (14.4,7.4) -- (14.4,6.0);
  \node[font=\scriptsize,color=grisnota,anchor=west] at (11.0,8.3) {escalón (red \net{STEP1} entre R47 y JP1)};
  % RSET y C1
  \draw (set) to[R,l=$R_{\mathrm{SET}}$,a={R4+R5+R6 = \SI{21}{\kilo\ohm}},*-] (14.4,3.0) -- (iepe);
  \draw (set) -- (11.0,1.4) to[C,l_={C1 \SI{100}{\nano\farad}}] (14.4,1.4) -- (14.4,3.0);
  % nodo IEPE, TP1, AFE y micrófono
  \draw (iepe) to[short,*-] (16.8,4.6) coordinate (bnc);
  \draw (iepe) node[below right,font=\scriptsize]{\net{IEPE}};
  \draw (15.2,4.6) to[short,*-o] (15.2,5.6) node[above,font=\scriptsize]{TP1};
  \draw (bnc) to[V,l_={$V_B+S\,p(t)$}] (16.8,1.6) node[ground]{};
  \node[font=\scriptsize,anchor=west] at (17.3,3.1) {micrófono (J2)};
  \draw[-{Stealth[length=2mm]},thick,color=verdeplaca] (16.8,4.6) -- (18.0,4.6) node[right,font=\scriptsize,align=left]{a R13\\(AFE)};
\end{circuitikz}
\caption{Fuente IEPE del canal 1. $R_{\mathrm{OUT}} = \mathrm{R2}\parallel\mathrm{R3} = \SI{75}{\ohm}$;
$R_{\mathrm{SET}}$ son tres resistencias en serie (\net{SET_I}–R4–\net{SET_MID}–R5–\net{SET_MID2}–R6–\net{IEPE}).
Con el puente JP1 puesto, R47 queda en paralelo con $R_{\mathrm{OUT}}$. El micrófono se dibuja con su modelo de la
ecuación~\eqref{eq:c03-iepe}.}
\label{fig:c03-fuente}
\end{figure}

\subsection{Cálculo de la corriente de cada canal}

Los valores salen de \archivo{tools/target.json}. En el canal 1, $R_{\mathrm{SET}}$ está hecha con tres resistencias en
serie y $R_{\mathrm{OUT}}$ con dos en paralelo (así viene del Altium del equipo). En el canal 2 cada una es una sola
pieza.

\begin{calculo}[corriente del canal 1 (RG-50)]
\begin{align*}
  R_{\mathrm{SET}} &= R_4 + R_5 + R_6 = \SI{10}{\kilo\ohm} + \SI{10}{\kilo\ohm} + \SI{1}{\kilo\ohm} = \SI{21}{\kilo\ohm}\\
  R_{\mathrm{OUT}} &= R_2 \parallel R_3 = \frac{150\cdot 150}{150+150}\,\si{\ohm} = \SI{75}{\ohm}\\
  I_{\mathrm{CH1}} &= \SI{10}{\micro\ampere}\left(1 + \frac{21000}{75}\right) = \SI{10}{\micro\ampere}\cdot 281
                    = \boxed{\SI{2.81}{\milli\ampere}}
\end{align*}
Con el puente JP1 puesto, R47 (\SI{180}{\ohm}) queda en paralelo con los \SI{75}{\ohm}:
\begin{align*}
  R_{\mathrm{OUT}}' &= \frac{75\cdot 180}{75+180}\,\si{\ohm} = \SI{52.94}{\ohm}\\
  I_{\mathrm{CH1}}' &= \SI{10}{\micro\ampere}\left(1 + \frac{21000}{52.94}\right) = \SI{10}{\micro\ampere}\cdot 397.7
                     = \boxed{\SI{3.98}{\milli\ampere}}
\end{align*}
\end{calculo}

\begin{calculo}[corriente del canal 2 (SV 12)]
En el canal 2, $R_{\mathrm{SET}}$ = R34 = \SI{15}{\kilo\ohm} (entre \net{SET_SV} e \net{IEPE_SV}) y
$R_{\mathrm{OUT}}$ = R35 = \SI{100}{\ohm} (entre \net{OUT_SV} e \net{IEPE_SV}):
\begin{align*}
  I_{\mathrm{CH2}} &= \SI{10}{\micro\ampere}\left(1 + \frac{15000}{100}\right) = \SI{10}{\micro\ampere}\cdot 151
                    = \boxed{\SI{1.51}{\milli\ampere}}
\end{align*}
Con JP2 puesto, R48 (\SI{150}{\ohm}) queda en paralelo con R35:
\begin{align*}
  R_{\mathrm{OUT}}' &= \frac{100\cdot 150}{100+150}\,\si{\ohm} = \SI{60}{\ohm}, \qquad
  I_{\mathrm{CH2}}' = \SI{10}{\micro\ampere}\left(1 + \frac{15000}{60}\right) = \boxed{\SI{2.51}{\milli\ampere}}
\end{align*}
\end{calculo}

Estas cuentas coinciden con las notas del esquemático (\archivo{tools/gen_sch_wired.py}) y con la tabla de entradas
de \archivo{CALCULOS.md}. La tolerancia de las resistencias es del \SI{1}{\percent} (son de la serie
\texttt{0603WAF…} y \texttt{RC0603FR…}); la de la corriente interna de \SI{10}{\micro\ampere} del LT3092 no está en los
textos de datasheet que se leyeron para este informe, así que el valor exacto de cada placa se mide con un
amperímetro en J1 o J7 \medir.

\subsection{Por qué esas corrientes}

\begin{decision}[la corriente de cada micrófono]
\textbf{Qué se decidió.} Canal 1 (ROGA RG-50): \SI{2.81}{\milli\ampere}. Canal 2 (Svantek SV 12): \SI{1.51}{\milli\ampere}.
Un puente en cada canal sube la corriente a \SI{3.98}{\milli\ampere} y \SI{2.51}{\milli\ampere}.

\textbf{Alternativas.} Una sola corriente para los dos (por ejemplo \SI{2.5}{\milli\ampere}, que documentan los dos
fabricantes); una corriente ajustable con potenciómetro.

\textbf{Por qué.}
\begin{itemize}
  \item El RG-50 acepta \SIrange{2}{20}{\milli\ampere} según su ficha actual, o \SIrange{2}{6}{\milli\ampere} según la
        versión vieja \doc{ROGA RG-50}. \SI{2.81}{\milli\ampere} y \SI{3.98}{\milli\ampere} caben en los dos rangos.
  \item El SV 12 se usa con \SI{1.5}{\milli\ampere} en el SVAN~977 y con \SI{2.5}{\milli\ampere} en el SVAN~958A
        \doc{Svantek}. Los dos puntos del canal 2 (\SI{1.51}{\milli\ampere} y \SI{2.51}{\milli\ampere}) son
        justamente esos dos valores documentados.
  \item El escalón (cambio C9 de la Rev B) sirve para la demostración: al cambiar la corriente cambia un poco el nivel
        de DC del nodo \net{IEPE}, y se puede ver en el osciloscopio cómo el lazo de DC (sección~\ref{sec:c03-lazo}) lo
        corrige. Se hizo con un puente de pines y no con un potenciómetro porque un puente da dos valores conocidos y
        repetibles.
\end{itemize}

\textbf{Fuente.} AGENTS.md §3 (fichas de ROGA y Svantek), \archivo{tools/design.py} (cambio C9).
\end{decision}

\subsection{La corriente que de verdad llega al micrófono}\label{sec:c03-carga}

El nodo \net{IEPE} no solo alimenta al micrófono: también ve la entrada del AFE, que es una cadena de
\SI{1.0009}{\mega\ohm} hasta el nodo \net{SUM}, que está a la tensión \net{VSHIFT} (el sumador es un inversor, su
entrada inversora es una tierra virtual a \net{VSHIFT}). Con el RG-50 ($V_B = \SI{12.5}{\volt}$) y
$V_{\mathrm{SHIFT}} \approx \SI{2.5}{\volt}$:
\begin{equation*}
  I_{\mathrm{AFE}} = \frac{V_B - V_{\mathrm{SHIFT}}}{R_{\mathrm{in}}} = \frac{\SI{12.5}{\volt}-\SI{2.5}{\volt}}{\SI{1.0009}{\mega\ohm}}
  \approx \SI{10}{\micro\ampere}.
\end{equation*}
Es decir, el micrófono recibe unos \SI{2.80}{\milli\ampere} en vez de \SI{2.81}{\milli\ampere}. Es el mismo número que
anota \archivo{docs/connection_matrix.md} (``\SI{2.810}{\milli\ampere} source; $\sim$\SI{2.800}{\milli\ampere}
sensor''). Esta es una de las razones de que la entrada del AFE sea de \SI{1}{\mega\ohm}: casi no le roba corriente
al micrófono.

\subsection{El resto de piezas de la fuente}

\begin{itemize}
  \item \textbf{J1 (bornera de 2 polos).} Está en serie entre \net{24V_PROT} y \net{IEPE_FEED}: con un puente de cable
        el canal funciona; quitándolo se apaga; y en su lugar se puede poner el amperímetro para medir la corriente.
        En el canal 2 es J7.
  \item \textbf{R1 (\SI{10}{\ohm}) con C2 (\SI{1}{\micro\farad}) y C3 (\SI{100}{\nano\farad}).} Forman un filtro RC en
        la entrada del LT3092 (red \net{IEPE_SUP}). Su corte es
        $f_c = 1/(2\pi\cdot\SI{10}{\ohm}\cdot\SI{1.1}{\micro\farad}) \approx \SI{14.5}{\kilo\hertz}$. El datasheet pide
        un bypass de entrada de \SIrange{0.1}{1}{\micro\farad} \copia[REGLAS.md §3]. El LT3092 no necesita
        condensadores para ser estable.
  \item \textbf{C1 (\SI{100}{\nano\farad}) en paralelo con $R_{\mathrm{SET}}$.} Filtra el ruido de la referencia de
        \SI{10}{\micro\ampere}. Una simulación previa del repositorio con el modelo de ADI
        (\archivo{docs/simulacion-lt3092.md}, marcada como provisional) dio \SI{3.99}{\micro\volt} rms de ruido en el
        nodo IEPE sin condensador y \SI{0.46}{\micro\volt} rms con \SI{100}{\nano\farad}, entre \SI{20}{\hertz} y
        \SI{20}{\kilo\hertz} \estimado. El piso de ruido del micrófono en esa simulación era
        \SI{40.8}{\micro\volt} rms, así que con C1 la fuente aporta alrededor del \SI{1}{\percent}.
  \item \textbf{TP1.} Punto de prueba en \net{IEPE} para el osciloscopio (cambio C10). En el canal 2 es TP3.
  \item \textbf{J2 (BNC a \SI{90}{\degree}).} Centro = \net{IEPE}, carcasa = GND (cambio C8). El SV 12 tiene conector
        TNC y se conecta a J8 con un adaptador BNC–TNC.
\end{itemize}

\subsection{¿Alcanzan los 24 V? (compliance)}

La fuente de corriente solo funciona si tiene tensión suficiente entre su entrada y el micrófono. Si el micrófono
sube su tensión (en un pico de sonido fuerte) por encima de lo que la fuente puede dar, la corriente se cae y la señal
se recorta.

\begin{calculo}[margen de tensión con el RG-50]
\begin{itemize}
  \item Tensión que llega: \SI{24}{\volt} menos la caída del diodo D1, unos \SI{23.5}{\volt} (\archivo{CALCULOS.md}).
  \item Caída de la fuente: unos \SI{1.42}{\volt} (\SI{1.22}{\volt} del integrado y \SI{0.2}{\volt} en
        $R_{\mathrm{OUT}}$) según la simulación previa del repositorio \estimado. A eso se suman unos milivoltios en R1.
  \item Tensión máxima disponible en el nodo IEPE: $\approx \SI{23.5}{\volt} - \SI{1.42}{\volt} \approx \SI{22}{\volt}$.
  \item Lo que pide el RG-50 en el peor caso: $V_B + v_{\mathrm{pico}} = \SI{12.5}{\volt} + \SI{3.2}{\volt}
        = \SI{15.7}{\volt}$ (\SI{130}{\decibel} pico con \SI{50}{\milli\volt\per\pascal}, AGENTS.md §3.2).
  \item Margen: unos \SI{6}{\volt}.
\end{itemize}
\end{calculo}

\begin{ojo}[el SV 12 todavía no está verificado]
Para el SV 12 Svantek publica una excursión de al menos \SI{10.8}{\volt} pico con \SIrange{28}{30}{\volt} de
alimentación, pero no publica su tensión de polarización. Con unos \SI{22}{\volt} disponibles, el SV 12 aprovecharía
toda su excursión solo si su polarización es menor que unos \SI{11}{\volt}. Hay que medir la polarización del SV 12 a
\SI{2.5}{\milli\ampere} \medir. Para el uso normal del sonómetro (niveles muy por debajo del máximo del micrófono) esto
no es un problema.
\end{ojo}

La disipación del LT3092 en el peor caso (micrófono en cortocircuito: toda la tensión cae en el integrado) es
$\SI{3.98}{\milli\ampere}\times\SI{23.3}{\volt} = \SI{93}{\milli\watt}$, que con $\theta_{JA} = \SI{57}{\degreeCelsius\per\watt}$
\copia{} sube la temperatura unos \SI{5.3}{\degreeCelsius} (\archivo{CALCULOS.md}, ver
capítulo~\ref{cap:calculos}).

% =====================================================================================================
\section{El sumador inversor: restar la polarización}\label{sec:c03-sumador}

\subsection{Dos maneras de quitar el DC}

Hay dos caminos clásicos para quedarte solo con $S\cdot p(t)$:
\begin{enumerate}
  \item \textbf{Acoplamiento por condensador.} Un condensador en serie bloquea el DC y deja pasar la señal; luego la
        recentras en la mitad del rango del ADC con un divisor. Es lo que hacía la propuesta v1 del grupo.
  \item \textbf{Acoplamiento en DC con resta de una referencia.} Un amplificador resta una tensión conocida y la
        señal pasa sin condensadores.
\end{enumerate}

El profesor criticó la v1 por dos cosas (AGENTS.md §10): el condensador de acoplamiento introduce desfase a baja
frecuencia, y el nivel de DC a mitad de escala ``no tenía una referencia definida''. Su sugerencia fue un
amplificador que reste (diferencial) para rechazar ruido.

El circuito del equipo, que es el que se ruteó, usa el segundo camino: \textbf{el sumador U3A resta la referencia
\net{VSHIFT} en DC}. Esa referencia la fija un DAC, el firmware sabe exactamente cuánto vale y además se mide
(\net{DC_MON}). Así el nivel de DC queda definido y controlado.

\begin{ojo}[sí queda un condensador, pero muy abajo en frecuencia]
Entre \net{REST_OUT} y el ADC está C9 (\SI{10}{\micro\farad}), porque el PCM1808 polariza su propia entrada a
$0{,}5\,V_{CC}$ y tiene una entrada de \SI{60}{\kilo\ohm} \doc{SLES177B p.6}. El corte de ese acoplamiento es
$1/(2\pi\cdot\SI{10}{\micro\farad}\cdot\SI{60}{\kilo\ohm}) \approx \SI{0.27}{\hertz}$, y el propio PCM1808 tiene un
filtro pasa altos digital de \SI{0.91}{\hertz} \doc{SLES177B p.1}. Los dos están dos décadas por debajo de
\SI{20}{\hertz}. El detalle del ADC está en el capítulo~\ref{cap:digital}.
\end{ojo}

\subsection{El circuito}

La figura~\ref{fig:c03-afe} es el AFE completo del canal 1. Fíjate en los nombres de las redes, porque son los que
vas a ver en la placa.

\begin{figure}[htbp]
\centering
\begin{circuitikz}[scale=0.78,transform shape,font=\small]
  % cadena de entrada
  \draw (0,0) node[left]{\net{IEPE}} to[R,l=R13,a=\SI{390}{\kilo\ohm},o-] (2.2,0)
        to[R,l=R14,a=\SI{365}{\kilo\ohm}] (4.4,0)
        to[R,l=R16,a=\SI{191}{\kilo\ohm}] (6.6,0)
        to[R,l=R18,a=\SI{54.9}{\kilo\ohm}] (8.8,0) coordinate (sum);
  \node[font=\scriptsize,below] at (2.2,-0.6) {\net{IEPE_DIV1}};
  \node[font=\scriptsize,below] at (4.4,-0.6) {\net{IEPE_DIV2}};
  \node[font=\scriptsize,below] at (6.6,-0.6) {\net{IEPE_DIV3}};
  \draw (2.2,0) node[circ]{}; \draw (4.4,0) node[circ]{}; \draw (6.6,0) node[circ]{};
  % U3A
  \draw (10.6,-0.48) node[op amp] (A) {\scriptsize U3A};
  \draw (sum) to[short,*-] (A.-);
  \draw (sum) node[above left,font=\scriptsize]{\net{SUM}};
  % realimentación
  \draw (sum) -- (8.8,2.0) to[R,l=R7,a=\SI{49.9}{\kilo\ohm}] (10.8,2.0) coordinate (fbm)
        to[R,l=R8,a=\SI{49.9}{\kilo\ohm}] (12.8,2.0) -- (12.8,-0.48);
  \draw (fbm) node[circ]{} node[above,font=\scriptsize]{\net{FB_MID}};
  \draw (8.8,2.0) -- (8.8,3.4) to[C,l={C56, \SI{10}{\pico\farad} C0G}] (12.8,3.4) -- (12.8,2.0);
  \draw (A.out) -- (12.8,-0.48) to[short,*-] (14.6,-0.48) coordinate (rest);
  \draw (12.8,-0.48) node[below right,font=\scriptsize]{\net{REST_OUT}};
  % salidas de REST_OUT
  \draw (13.7,-0.48) to[short,*-o] (13.7,0.6) node[above,font=\scriptsize]{TP2};
  \draw (rest) to[C,l=C9,a=\SI{10}{\micro\farad}] (16.6,-0.48) node[right,font=\scriptsize,align=left]{\net{AC_LEFT}\\(PCM1808)};
  \draw (14.6,-0.48) -- (14.6,-1.6) to[R,l_={R15 \SI{10}{\kilo\ohm}}] (14.6,-3.6) coordinate (mon)
        to[R,l_={R17 \SI{10}{\kilo\ohm}}] (14.6,-5.6) node[ground]{};
  \draw (mon) to[short,*-] (16.2,-3.6) to[C,l={C8 \SI{100}{\nano\farad}}] (16.2,-5.6) node[ground]{};
  \draw (16.2,-3.6) -- (17.2,-3.6) node[right,font=\scriptsize,align=left]{\net{DC_MON}\\NUCLEO A0};
  % + de U3A: VSHIFT desde el seguidor U3B
  \draw (A.+) -- ++(-0.5,0) coordinate (vsh);
  \draw (6.2,-3.0) node[op amp,yscale=-1] (B) {};
  \node[font=\scriptsize] at (B.center) {U3B};
  \draw (B.out) -| (vsh);
  \draw (B.out) to[short,*-] ++(0,-1.25) coordinate (bfb) -| (B.-);
  \node[font=\scriptsize,above] at (8.1,-3.0) {\net{VSHIFT}};
  % DAC y filtro
  \draw (B.+) -- (3.6,-2.52) coordinate (dacf);
  \draw (dacf) node[circ]{} node[above,font=\scriptsize]{\net{DAC_F}};
  \draw (dacf) -- (3.6,-4.6) to[C,l=C4,a=\SI{100}{\nano\farad}] (3.6,-6.4) node[ground]{};
  \draw (dacf) to[R,l_=R12,a^=\SI{10}{\kilo\ohm}] (1.0,-2.52) node[left,font=\scriptsize,align=right]{\net{DAC_RAW}\\(U2, MCP4725)};
\end{circuitikz}
\caption{AFE del canal 1 con las redes reales. $R_{\mathrm{in}}$ = R13+R14+R16+R18 = \SI{1000.9}{\kilo\ohm};
$R_f$ = R7+R8 = \SI{99.8}{\kilo\ohm}, en paralelo con C56. U3A y U3B son las dos mitades del mismo TLV9152 (U3),
alimentado de \net{5V_A} (pin 8) y GND (pin 4), desacoplado con C7. Pines: U3.2 = \net{SUM}, U3.3 = \net{VSHIFT},
U3.1 = \net{REST_OUT}, U3.5 = \net{DAC_F}, U3.6 = U3.7 = \net{VSHIFT}.}
\label{fig:c03-afe}
\end{figure}

Componentes del sumador:
\begin{itemize}
  \item \textbf{Resistencia de entrada} $R_{\mathrm{in}}$ = R13 + R14 + R16 + R18
        = \SI{390}{\kilo\ohm} + \SI{365}{\kilo\ohm} + \SI{191}{\kilo\ohm} + \SI{54.9}{\kilo\ohm} = \SI{1000.9}{\kilo\ohm}.
  \item \textbf{Resistencia de realimentación} $R_f$ = R7 + R8 = 2 × \SI{49.9}{\kilo\ohm} = \SI{99.8}{\kilo\ohm}.
  \item \textbf{CF} = C56, \SI{10}{\pico\farad} C0G, en paralelo con $R_f$ (agregado en la Rev B, sección~\ref{sec:c03-cf}).
  \item \textbf{U3A}: mitad A del TLV9152. Entrada $-$ en \net{SUM}, entrada $+$ en \net{VSHIFT}, salida \net{REST_OUT}.
\end{itemize}

Que $R_{\mathrm{in}}$ esté partida en cuatro resistencias y $R_f$ en dos viene del diseño del equipo; el repositorio no
documenta el motivo. Lo que sí importa para el ruteo es que los nodos intermedios (\net{IEPE_DIV1}…\net{IEPE_DIV3},
\net{SUM}) son de alta impedancia, así que son sensibles a acoples. Por eso se rutearon solo en \capa{F.Cu} y nada
pasa por debajo de ellos (capítulo~\ref{cap:ruteo}).

\subsection{La ecuación, paso a paso}

Vamos a deducir \net{REST_OUT} con superposición, como en el curso de circuitos. Suponemos el op-amp ideal: ganancia
infinita, entradas que no consumen corriente, y por eso las dos entradas a la misma tensión. Llamemos
$k = R_f/R_{\mathrm{in}}$.

\begin{calculo}[superposición en el sumador]
\textbf{Paso 1: solo \net{IEPE}} (\net{VSHIFT} a 0~V). Es un inversor clásico:
\begin{equation*}
  V_{\mathrm{REST}}^{(1)} = -\frac{R_f}{R_{\mathrm{in}}}\,v_{\mathrm{IEPE}} = -k\,v_{\mathrm{IEPE}}.
\end{equation*}
\textbf{Paso 2: solo \net{VSHIFT}} (\net{IEPE} a 0~V). Visto desde la entrada $+$ es un no inversor:
\begin{equation*}
  V_{\mathrm{REST}}^{(2)} = \left(1 + \frac{R_f}{R_{\mathrm{in}}}\right) V_{\mathrm{SHIFT}} = (1+k)\,V_{\mathrm{SHIFT}}.
\end{equation*}
\textbf{Paso 3: se suman}:
\begin{equation}
  \boxed{V_{\mathrm{REST}} = (1+k)\,V_{\mathrm{SHIFT}} - k\,v_{\mathrm{IEPE}}
  = V_{\mathrm{SHIFT}} - k\left(v_{\mathrm{IEPE}} - V_{\mathrm{SHIFT}}\right)}
  \label{eq:c03-rest}
\end{equation}
\textbf{Paso 4: los números}:
\begin{equation*}
  k = \frac{\SI{99.8}{\kilo\ohm}}{\SI{1000.9}{\kilo\ohm}} = \num{0.0997}, \qquad 1+k = \num{1.0997}.
\end{equation*}
\end{calculo}

Ahora mete la ecuación~\eqref{eq:c03-iepe} en la~\eqref{eq:c03-rest}:
\begin{equation*}
  V_{\mathrm{REST}} = \underbrace{(1+k)\,V_{\mathrm{SHIFT}} - k\,V_B}_{\text{nivel de DC}}
  \;\underbrace{-\;k\,S\,p(t)}_{\text{señal}}.
\end{equation*}
Esto dice tres cosas:
\begin{itemize}
  \item La \textbf{señal} sale multiplicada por $-k \approx -0{,}1$: atenuada diez veces e invertida (la inversión no
        importa, el firmware la puede deshacer).
  \item El \textbf{nivel de DC} depende de $V_{\mathrm{SHIFT}}$ y de $V_B$. Como $V_B$ depende del micrófono, de la
        corriente y de la temperatura, no se puede fijar de una vez: por eso hay un lazo que ajusta
        $V_{\mathrm{SHIFT}}$ (sección~\ref{sec:c03-lazo}).
  \item Para que la salida quede en un objetivo $V_{\mathrm{REST}}^{*}$ hay que poner
        \begin{equation}
          V_{\mathrm{SHIFT}} = \frac{V_{\mathrm{REST}}^{*} + k\,V_B}{1+k}.
          \label{eq:c03-vshift}
        \end{equation}
\end{itemize}

\begin{calculo}[ejemplo con el RG-50]
$V_B = \SI{12.5}{\volt}$ \doc{ROGA RG-50} y un objetivo de \SI{1.5}{\volt} en \net{REST_OUT}:
\begin{equation*}
  V_{\mathrm{SHIFT}} = \frac{1{,}5 + 0{,}0997\cdot 12{,}5}{1{,}0997}\,\si{\volt} = \frac{1{,}5 + 1{,}246}{1{,}0997}\,\si{\volt}
  = \SI{2.50}{\volt}.
\end{equation*}
Con \SI{94}{\decibel} SPL (\SI{1}{\pascal} rms, el nivel típico de un calibrador) y \SI{50}{\milli\volt\per\pascal},
el micrófono da \SI{50}{\milli\volt} rms y en \net{REST_OUT} quedan
$0{,}0997\times\SI{50}{\milli\volt} \approx \SI{5}{\milli\volt}$ rms sobre \SI{1.5}{\volt}.
\end{calculo}

\subsection{Por qué la ganancia es 0,1}

\begin{decision}[atenuar diez veces]
\textbf{Qué se decidió.} $k = R_f/R_{\mathrm{in}} \approx 0{,}1$ (diseño del equipo, conservado en la Rev B).

\textbf{Alternativas.} Ganancia 1 o mayor (como la etapa de G = 11 de la propuesta v1, que el profesor criticó).

\textbf{Por qué.} La ganancia se elige para que \emph{el peor caso quepa} en el ADC:
\begin{itemize}
  \item El PCM1808 acepta \SI{3}{\volt} pico a pico a fondo de escala, es decir $\pm\SI{1.5}{\volt}$
        \doc{SLES177B p.4, p.6}.
  \item El SV 12 puede excursionar al menos \SI{10.8}{\volt} pico \doc{Svantek}. Por 0,1 quedan
        \SI{1.08}{\volt} pico: cabe con margen.
  \item El RG-50 a su máximo de \SI{130}{\decibel} pico da unos \SI{3.2}{\volt} pico; por 0,1 quedan
        \SI{0.32}{\volt}.
\end{itemize}
Con ganancia 1 la excursión del SV 12 recortaría el ADC diez veces antes.

\textbf{Fuente.} AGENTS.md §3, SLES177B, \archivo{tools/analisis_opamp.py} (ganancia $-0{,}0997$).
\end{decision}

\begin{ojo}[lo que se paga por atenuar]
Atenuar diez veces también acerca la señal al ruido del ADC. Una estimación gruesa \estimado: el PCM1808 tiene
\SI{99}{\decibel} de rango dinámico típico \doc{SLES177B p.1}; sobre \SI{3}{\volt} pico a pico
(\SI{1.06}{\volt} rms) eso es un piso de unos \SI{12}{\micro\volt} rms. Llevado a la entrada del RG-50
($0{,}0997\times\SI{50}{\milli\volt\per\pascal}$) equivale a unos \SI{2.4}{\milli\pascal}, unos \SI{41}{\decibel} SPL. El
ruido propio del RG-50 es de unos \SI{27}{\decibel}(A). Para niveles normales de sala y de calibrador (\SI{94}{\decibel})
sobra, pero el piso del sistema lo pone el ADC y no el micrófono. Son números típicos y con ponderaciones distintas: hay
que medirlo \medir.
\end{ojo}

\subsection{Qué pasa si se desconecta el micrófono}

Si el cable se abre, la fuente de corriente ya no tiene a dónde mandar su corriente y el nodo \net{IEPE} sube casi
hasta la entrada (\
um{22}–\SI{24}{\volt}). Con la ecuación~\eqref{eq:c03-rest}, \net{REST_OUT} pediría una tensión
negativa: el op-amp se satura contra GND. Lo bueno es que la cadena de \SI{1}{\mega\ohm} protege la entrada del
op-amp: aun en el peor caso, la corriente hacia \net{SUM} es de unos $\SI{24}{\volt}/\SI{1}{\mega\ohm} =
\SI{24}{\micro\ampere}$, muy por debajo de los \SI{10}{\milli\ampere} que el TLV9152 admite en sus entradas como máximo
absoluto \doc{SBOS986E p.9}. Al reconectar el micrófono, el lazo de DC vuelve a centrar la salida.

% =====================================================================================================
\section{La referencia VSHIFT: DAC, filtro y seguidor}\label{sec:c03-vshift}

\subsection{El DAC MCP4725}

U2 (canal 1) y U10 (canal 2) son MCP4725: DAC de 12 bits con salida de tensión, interfaz I\textsuperscript{2}C y
memoria EEPROM \doc{DS22039D p.1}. Su referencia es su propia alimentación ($V_{\mathrm{REF}} = V_{DD}$)
\doc{DS22039D p.15}, así que
\begin{equation*}
  V_{\mathrm{DAC}} = V_{DD}\cdot\frac{\text{código}}{4096}, \qquad
  1\ \text{LSB} = \frac{\SI{3.3}{\volt}}{4096} \approx \SI{0.81}{\milli\volt}.
\end{equation*}
Como \net{REST_OUT} cambia $(1+k)$ veces lo que cambia \net{VSHIFT}, un LSB del DAC mueve la salida unos
$1{,}1\times\SI{0.81}{\milli\volt} \approx \SI{0.89}{\milli\volt}$. Esa es la resolución con la que el lazo puede
centrar el DC.

\begin{figure}[htbp]
\centering
\begin{circuitikz}[scale=0.9,transform shape,font=\small]
  \draw (0,3) node[left]{\net{3V3}} to[R,l=R9,a=\SI{100}{\ohm},o-] (2.6,3) coordinate (vdd);
  \draw (vdd) node[above,font=\scriptsize]{\net{3V3_DAC}};
  \draw (vdd) to[short,*-] (2.6,2.4) to[C,l_={C5 \SI{10}{\micro\farad}}] (2.6,0.4) node[ground]{};
  \draw (vdd) -- (4.0,3) to[short,*-] (4.0,2.4) to[C,l={C6 \SI{100}{\nano\farad}}] (4.0,0.4) node[ground]{};
  \draw (4.0,3) -- (6.0,3);
  \draw[thick] (6.0,0.2) rectangle (8.6,3.6);
  \node[align=center] at (7.3,1.9) {U2\\MCP4725};
  \node[font=\scriptsize,right] at (6.0,3.0) {3 VDD};
  \node[font=\scriptsize,right] at (6.0,0.6) {2 VSS};
  \draw (6.0,0.6) -- (5.4,0.6) node[ground]{};
  \node[font=\scriptsize,left] at (8.6,3.0) {VOUT 1};
  \node[font=\scriptsize,left] at (8.6,2.2) {SDA 4};
  \node[font=\scriptsize,left] at (8.6,1.4) {SCL 5};
  \node[font=\scriptsize,left] at (8.6,0.6) {A0 6};
  \draw (8.6,3.0) -- (9.4,3.0) node[right,font=\scriptsize]{\net{DAC_RAW} $\to$ R12};
  \draw (8.6,0.6) -- (9.4,0.6) node[ground]{};
  \node[font=\scriptsize,right] at (9.5,0.65) {(CH2: A0 a \net{3V3_DAC_SV} $\Rightarrow$ 0x61)};
  \draw (8.6,2.2) -- (12.6,2.2) coordinate (sda);
  \draw (8.6,1.4) -- (14.0,1.4) coordinate (scl);
  \draw (sda) to[R,l_=R11,a^=\SI{4.7}{\kilo\ohm},*-] (12.6,4.8) node[above,font=\scriptsize]{\net{3V3}};
  \draw (scl) to[R,l=R10,a=\SI{4.7}{\kilo\ohm},*-] (14.0,4.8) node[above,font=\scriptsize]{\net{3V3}};
  \draw (sda) -- (15.4,2.2) node[right,font=\scriptsize]{\net{I2C_SDA}};
  \draw (scl) -- (15.4,1.4) node[right,font=\scriptsize]{\net{I2C_SCL}};
\end{circuitikz}
\caption{DAC del canal 1: alimentación filtrada (R9, C5, C6), dirección por A0 y pull-ups del bus I\textsuperscript{2}C.
El bus es compartido: U10 (canal 2) cuelga de las mismas \net{I2C_SDA}/\net{I2C_SCL}. En el canal 2 la alimentación es
R46, C52 y C53.}
\label{fig:c03-dac}
\end{figure}

\paragraph{La dirección I\textsuperscript{2}C.} La dirección del MCP4725 es \texttt{1100} + A2 A1 A0. A2 y A1 los graba
la fábrica (\texttt{00} en la variante A0 que se compró, \net{MCP4725A0T-E/CH}) y A0 lo fija el pin 6
\doc{DS22039D p.27}. En U2 el pin 6 va a GND: dirección \texttt{1100000} = \textbf{0x60}. En U10 va a
\net{3V3_DAC_SV}: \texttt{1100001} = \textbf{0x61}. Así los dos DAC comparten el bus sin chocar.

\paragraph{La EEPROM.} El MCP4725 guarda un código en su memoria no volátil y lo carga al encender
\doc{DS22039D p.1}. Si el firmware guarda ahí el último \net{VSHIFT} bueno, el canal arranca ya cerca de su punto de
trabajo. (Escribir la EEPROM tarda hasta \SI{50}{\milli\second} \doc{DS22039D p.4}, así que no se hace en cada
corrección del lazo.)

\paragraph{La alimentación del DAC es su referencia.} Cualquier ruido en $V_{DD}$ aparece en la salida del DAC
\doc{DS22039D §8.3}, y de ahí pasa a \net{REST_OUT} multiplicado por 1,1. Por eso el DAC no se alimenta directo del
riel de \SI{3.3}{\volt} del buck, sino a través de R9 (\SI{100}{\ohm}) con C5 (\SI{10}{\micro\farad}) y C6
(\SI{100}{\nano\farad}):
\begin{equation*}
  f_c = \frac{1}{2\pi\cdot\SI{100}{\ohm}\cdot\SI{10.1}{\micro\farad}} \approx \SI{158}{\hertz}.
\end{equation*}
El buck conmuta a \SI{1}{\mega\hertz} (capítulo~\ref{cap:alim}); a esa frecuencia el filtro atenúa unas
$\num{6300}$ veces (más de \SI{70}{\decibel}) en la aproximación de primer orden. El precio es una pequeña caída:
con \SI{210}{\micro\ampere} típicos de consumo \doc{DS22039D p.3}, R9 baja \SI{21}{\milli\volt}, así que
\net{3V3_DAC} queda en unos \SI{3.28}{\volt}. No importa: el lazo de DC mide el resultado y corrige.

\subsection{El filtro R12/C4}

La salida del DAC (\net{DAC_RAW}) pasa por R12 (\SI{10}{\kilo\ohm}) y C4 (\SI{100}{\nano\farad}) antes de llegar al
seguidor (\net{DAC_F}):
\begin{equation*}
  \tau = R_{12}\,C_4 = \SI{10}{\kilo\ohm}\cdot\SI{100}{\nano\farad} = \SI{1}{\milli\second}, \qquad
  f_c = \frac{1}{2\pi\tau} \approx \SI{159}{\hertz}.
\end{equation*}
Sirve para dos cosas: suaviza los saltos cuando el DAC cambia de código (el lazo no mete escalones bruscos en la señal)
y quita ruido de alta frecuencia del DAC.

\begin{decision}[filtros del lazo de \SI{10}{\micro\farad} a \SI{100}{\nano\farad} (cambio C6)]
\textbf{Qué se decidió.} C4/C51 (filtro del DAC) y C8/C55 (filtro de \net{DC_MON}) pasaron de \SI{10}{\micro\farad} a
\SI{100}{\nano\farad} X7R.

\textbf{Alternativas.} Dejar los \SI{10}{\micro\farad} del Altium original.

\textbf{Por qué.} Con \SI{10}{\micro\farad} las constantes de tiempo eran \SI{100}{\milli\second} (DAC) y
\SI{50}{\milli\second} (\net{DC_MON}). La corrección del DC tardaría cientos de milisegundos en asentarse. Con
\SI{100}{\nano\farad} quedan en \SI{1}{\milli\second} y \SI{0.5}{\milli\second}, y el lazo corrige en milisegundos.
Además, un \SI{100}{\nano\farad} X7R en 0603 es una pieza \emph{basic} de JLC (C14663).

\textbf{Fuente.} \archivo{tools/changes.json}, cambio C6.
\end{decision}

\subsection{El seguidor U3B}

La otra mitad del TLV9152 (pines 5, 6 y 7) está conectada como seguidor de tensión: la entrada $+$ (U3.5) recibe
\net{DAC_F}, y la entrada $-$ (U3.6) está unida a la salida (U3.7). La salida es \net{VSHIFT}, que va solo a la entrada
$+$ del sumador (U3.3). Puedes comprobarlo en \archivo{tools/target.json}: la red \net{VSHIFT} tiene exactamente tres
pines, U3.3, U3.6 y U3.7, y ningún condensador.

¿Para qué sirve, si la entrada $+$ de U3A casi no consume corriente? El repositorio no lo explica (viene del diseño del
equipo). Una ventaja práctica es que \net{VSHIFT} queda a baja impedancia: una punta de osciloscopio, una fuga en la
placa o la corriente de entrada del sumador no mueven su nivel, cosa que sí podría pasar en el nodo \net{DAC_F}, que
tiene \SI{10}{\kilo\ohm} de impedancia. Y como el TLV9152 es doble, esa mitad está disponible sin costo.

La estabilidad del seguidor se revisa en la sección~\ref{sec:c03-tlv}.

% =====================================================================================================
\section{DC\_MON y el lazo de DC}\label{sec:c03-lazo}

\subsection{Medir el nivel de DC}

R15 y R17 (\SI{10}{\kilo\ohm} cada una) dividen \net{REST_OUT} entre 2 y C8 (\SI{100}{\nano\farad}) filtra el
resultado. La red se llama \net{DC_MON} y llega al pin A0 (PA0) de la NUCLEO (capítulo~\ref{cap:digital}).
\begin{itemize}
  \item \textbf{Por qué dividir entre 2.} \net{REST_OUT} puede llegar hasta casi \SI{5}{\volt} (el op-amp va de
        \net{5V_A}); el ADC del STM32 trabaja de \SIrange{0}{3.3}{\volt}. Dividido, el máximo queda en unos
        \SI{2.5}{\volt}.
  \item \textbf{El filtro.} La resistencia equivalente del divisor es $R_{15}\parallel R_{17} = \SI{5}{\kilo\ohm}$; con C8,
        $\tau = \SI{0.5}{\milli\second}$ y $f_c \approx \SI{318}{\hertz}$. Atenúa la señal de audio pero no la elimina
        del todo: por eso el firmware promedia muchas muestras para quedarse con el DC.
  \item \textbf{Precisión.} R15/R17 se cambiaron de \SI{5}{\percent} a \SI{1}{\percent} al pasarlas a piezas
        \emph{basic} de JLC (C25804), así que la relación del divisor queda al \SI{1}{\percent} o mejor.
\end{itemize}

\subsection{El lazo en el firmware}

La nota del esquemático (\archivo{tools/gen_sch_wired.py}) dice: \emph{``Lazo de DC (firmware): DC\_MON $\to$ nuevo
VSHIFT = VSHIFT $-$ error/1.1 $\to$ DAC; quieto al medir''}. Así se entiende:

\begin{calculo}[un paso del lazo]
\begin{enumerate}
  \item El firmware mide \net{DC_MON} (promediado) y calcula $V_{\mathrm{REST}} = 2\cdot V_{\mathrm{DC\_MON}}$.
  \item Calcula el error contra el objetivo: $e = V_{\mathrm{REST}} - V_{\mathrm{REST}}^{*}$.
  \item De la ecuación~\eqref{eq:c03-rest}, \net{REST_OUT} cambia $(1+k) \approx 1{,}1$ veces lo que cambia
        \net{VSHIFT}. Para quitar todo el error en un solo paso:
        \begin{equation*}
          V_{\mathrm{SHIFT}}^{\text{nuevo}} = V_{\mathrm{SHIFT}} - \frac{e}{1+k} \approx V_{\mathrm{SHIFT}} - \frac{e}{1{,}1}.
        \end{equation*}
  \item Convierte a código del DAC ($V/V_{\mathrm{LSB}}$) y lo envía por I\textsuperscript{2}C.
\end{enumerate}
Ejemplo: objetivo \SI{1.5}{\volt}, se mide \SI{1.6}{\volt}. Error \SI{0.1}{\volt}. \net{VSHIFT} baja
$\SI{0.1}{\volt}/1{,}1 = \SI{91}{\milli\volt}$, unos 113 códigos del DAC.
\end{calculo}

``Quieto al medir'' significa que el lazo \emph{no} se mueve mientras se graba audio: cada cambio de \net{VSHIFT}
aparece en la señal (multiplicado por 1,1). El lazo corrige entre medidas, o cuando el DC se va demasiado (por
ejemplo al poner el puente JP1).

\begin{idea}[por qué esto responde a la crítica del profesor]
En la v1 el nivel de DC salía de un divisor de la alimentación y nadie lo medía. Aquí el nivel de DC es una consigna
del firmware: se elige, se genera con un DAC y se comprueba con un ADC. Eso es una ``referencia definida''. Y el canal
no tiene condensadores en el camino de la señal antes del ADC.
\end{idea}

El objetivo de \net{REST_OUT} no está fijado en el repositorio. La sección~\ref{sec:c03-transicion} explica por qué se
recomienda \SI{1.5}{\volt} y no la mitad del riel.

% =====================================================================================================
\section{Estabilidad del sumador: por qué hizo falta CF}\label{sec:c03-cf}

Esta es la parte más ``de control'' del capítulo. Si nunca viste margen de fase en un op-amp, lee primero la caja.

\begin{idea}[margen de fase con un op-amp]
Un op-amp realimentado es un sistema con lazo cerrado. La ganancia que da la vuelta al lazo es $A\beta$: $A$ es la
ganancia en lazo abierto del op-amp y $\beta$ la fracción de la salida que vuelve a la entrada inversora. Donde
$|A\beta| = 1$ (el \emph{cruce}), mira cuánta fase le falta a $A\beta$ para llegar a $-\SI{180}{\degree}$: eso es el
\textbf{margen de fase}. Con margen de 60–\SI{70}{\degree} el circuito responde limpio. Con pocos grados, la
respuesta tiene un pico grande y oscila al menor golpe; con 0 oscila solo.

Una forma práctica de verlo en el diagrama de Bode es dibujar $|A|$ y $|1/\beta|$ (la ``ganancia de ruido''). Se
cruzan donde $|A\beta| = 1$. Si en el cruce $|A|$ baja a \SI{-20}{\decibel}/década y $|1/\beta|$ está plana, el margen
es bueno. Si $|1/\beta|$ \emph{sube} a $+$\SI{20}{\decibel}/década en el cruce (``pendiente de cierre'' de
\SI{40}{\decibel}/década), el margen es malo.
\end{idea}

\subsection{El polo escondido en el nodo SUM}

El nodo \net{SUM} no es ideal: tiene capacidad a tierra. Hay capacidad dentro del op-amp (diferencial y de modo común)
y capacidad de los pads y la pista. En el modelo de \archivo{tools/analisis_opamp.py}:
\begin{equation*}
  C_{\mathrm{SUM}} = \underbrace{\SI{3}{\pico\farad}}_{C_{ID}} + \underbrace{\SI{1}{\pico\farad}}_{C_{ICM}}
  + \underbrace{\SI{1.5}{\pico\farad}}_{\text{pista}} = \SI{5.5}{\pico\farad}
\end{equation*}
($C_{ID}$ y $C_{ICM}$ \doc{SBOS986E p.11}; la pista \estimado).

\SI{5.5}{\pico\farad} parece nada, pero está en un nodo de \SI{1}{\mega\ohm}. La ganancia de ruido es
\begin{equation*}
  \frac{1}{\beta} = 1 + \frac{Z_f}{Z_{\mathrm{in}}} = 1 + R_f\left(\frac{1}{R_{\mathrm{in}}} + s\,C_{\mathrm{SUM}}\right)
  = 1{,}0997 + s\,R_f\,C_{\mathrm{SUM}} \qquad (\text{sin CF}).
\end{equation*}
Tiene un cero en
\begin{equation*}
  f_z = \frac{1{,}0997}{2\pi\,R_f\,C_{\mathrm{SUM}}} = \frac{1{,}0997}{2\pi\cdot\SI{99.8}{\kilo\ohm}\cdot\SI{5.5}{\pico\farad}}
  \approx \SI{319}{\kilo\hertz}.
\end{equation*}
Por encima de \SI{319}{\kilo\hertz}, $|1/\beta|$ sube a $+$\SI{20}{\decibel}/década. Según el datasheet \doc{SBOS986E p.12}, el
TLV9152 tiene GBW = \SI{4.5}{\mega\hertz}. Entonces $|A| \approx \mathrm{GBW}/f$, y las dos curvas se
cruzan donde
\begin{equation*}
  \frac{\mathrm{GBW}}{f} = 2\pi f R_f C_{\mathrm{SUM}} \;\Rightarrow\;
  f = \sqrt{\frac{\mathrm{GBW}}{2\pi R_f C_{\mathrm{SUM}}}} = \sqrt{\frac{\num{4.5e6}}{2\pi\cdot\num{99.8e3}\cdot\num{5.5e-12}}}
  \approx \SI{1.14}{\mega\hertz}.
\end{equation*}
Ahí la pendiente de cierre es \SI{40}{\decibel}/década: margen casi nulo. El script, que además incluye el segundo
polo del op-amp, da \SI{7.8}{\degree} de margen y un pico de \SI{17}{\decibel} (tabla~\ref{tab:c03-cf}).

\subsection{Cómo lo arregla CF}

Con C56 en paralelo con $R_f$, a alta frecuencia la realimentación es un divisor capacitivo:
\begin{equation*}
  \frac{1}{\beta}\Big|_{f\to\infty} \approx \frac{C_{\mathrm{SUM}} + C_F}{C_F} = \frac{5{,}5 + 10}{10} = 1{,}55.
\end{equation*}
La curva $|1/\beta|$ deja de subir y se queda plana en 1,55 (\SI{3.8}{\decibel}). Ahora cruza a $|A|$ cuando $|A|$ baja
a \SI{-20}{\decibel}/década y $|1/\beta|$ está plana: pendiente de cierre \SI{20}{\decibel}/década, margen bueno. El
cruce queda cerca de $\mathrm{GBW}/1{,}55 \approx \SI{2.9}{\mega\hertz}$ (el script da \SI{2.74}{\mega\hertz}) y el
margen sube a \SI{69.7}{\degree}. La figura~\ref{fig:c03-bode} muestra las tres curvas.

\begin{figure}[htbp]
\centering
\begin{tikzpicture}
\begin{axis}[
  width=0.92\linewidth, height=7.2cm,
  xmode=log, ymode=log,
  xmin=1e3, xmax=1e8, ymin=0.5, ymax=5000,
  xlabel={Frecuencia (Hz)}, ylabel={Magnitud (V/V)},
  grid=both, grid style={line width=.1pt, draw=gray!20}, major grid style={line width=.2pt,draw=gray!40},
  legend pos=north east, legend style={font=\scriptsize,fill=white,fill opacity=0.9,text opacity=1},
  tick label style={font=\scriptsize}, label style={font=\small}]
  % |A| = GBW/f (polo dominante; se omite el segundo polo)
  \addplot[thick,color=black,domain=3:8,samples=60] ({10^x},{4.5e6/(10^x)});
  \addlegendentry{$|A|$ TLV9152 ($\mathrm{GBW}/f$)}
  % |1/beta| sin CF
  \addplot[thick,color=rojoojo,domain=3:8,samples=200]
    ({10^x},{sqrt(1.0997^2 + (0.0997*2*pi*(10^x)*5.505e-6)^2)});
  \addlegendentry{$|1/\beta|$ sin CF}
  % |1/beta| con CF = 10 pF
  \addplot[thick,color=verdeplaca,domain=3:8,samples=200]
    ({10^x},{sqrt( (1 + 0.0997*(1 + (2*pi*(10^x)*5.505e-6)*(2*pi*(10^x)*9.98e-7))/(1+(2*pi*(10^x)*9.98e-7)^2))^2
                 + (0.0997*((2*pi*(10^x)*5.505e-6) - (2*pi*(10^x)*9.98e-7))/(1+(2*pi*(10^x)*9.98e-7)^2))^2 )});
  \addlegendentry{$|1/\beta|$ con CF = \SI{10}{\pico\farad}}
  \draw[densely dashed,color=rojoojo] (axis cs:1.14e6,0.5) -- (axis cs:1.14e6,3.95);
  \node[font=\scriptsize,color=rojoojo,anchor=east] at (axis cs:1.05e6,0.72) {cruce \SI{1.1}{\mega\hertz}, \SI{40}{\decibel}/déc.};
  \draw[densely dashed,color=verdeplaca] (axis cs:2.9e6,0.5) -- (axis cs:2.9e6,1.55);
  \node[font=\scriptsize,color=verdeplaca,anchor=west] at (axis cs:3.2e6,2.5) {cruce \SI{2.9}{\mega\hertz}};
\end{axis}
\end{tikzpicture}
\caption{Diagrama de Bode (solo magnitud) del lazo del sumador. Sin CF, la ganancia de ruido sube desde
\SI{319}{\kilo\hertz} y corta a $|A|$ con pendiente de cierre de \SI{40}{\decibel}/década. Con CF = \SI{10}{\pico\farad}
se aplana en 1,55 y el cruce es limpio. Curvas calculadas con los valores de la sección; $|A|$ ignora el segundo polo
del op-amp, que sí incluye \archivo{tools/analisis_opamp.py}.}
\label{fig:c03-bode}
\end{figure}

\subsection{Elegir el valor de CF}

El script barre varios valores de CF con un modelo de dos polos: el polo dominante sale del GBW y el segundo polo se
ajusta para reproducir el margen de fase a ganancia 1 del datasheet (\SI{60}{\degree} en el TLV9152 \doc{SBOS986E
p.12}, \SI{55}{\degree} en el TLV9062 \doc{SBOS839N p.13}). La tabla~\ref{tab:c03-cf} es su salida
(\texttt{python -X utf8 -I tools/analisis\_opamp.py}).

\begin{table}[htbp]
\centering\small
\caption{Barrido de CF en el sumador (salida de \archivo{tools/analisis_opamp.py}). ``$-3$~dB'' es el ancho de banda
de la señal de \net{IEPE} a \net{REST_OUT}.}
\label{tab:c03-cf}
\begin{tabular}{@{}r rrrr rrr@{}}
\toprule
 & \multicolumn{4}{c}{TLV9152 (el de la placa)} & \multicolumn{3}{c}{TLV9062 (el original)}\\
\cmidrule(lr){2-5}\cmidrule(l){6-8}
CF & cruce & margen & pico & $-3$ dB & cruce & margen & $-3$ dB\\
(\si{\pico\farad}) & (\si{\kilo\hertz}) & (\si{\degree}) & (\si{\decibel}) & (\si{\kilo\hertz}) & (\si{\kilo\hertz}) & (\si{\degree}) & (\si{\kilo\hertz})\\
\midrule
0    & 1114 & 7,8  & 17,0 & 1820 & 1445 & 3,4  & 2291\\
4,7  & 2027 & 70,8 & 0    & 380  & 3739 & 72,3 & 363\\
\textbf{10}   & \textbf{2741} & \textbf{69,7} & \textbf{0} & \textbf{166} & 5352 & 68,8 & 166\\
22   & 3313 & 66,8 & 0    & 75,9 & 6744 & 64,6 & 75,9\\
47   & 3649 & 64,9 & 0    & 34,7 & 7611 & 61,9 & 34,7\\
\bottomrule
\end{tabular}
\end{table}

\begin{figure}[htbp]
\centering
\begin{tikzpicture}
\begin{axis}[
  width=0.8\linewidth, height=5.6cm,
  xlabel={CF (\si{\pico\farad})}, ylabel={Margen de fase (\si{\degree})},
  xmin=-2, xmax=50, ymin=0, ymax=80,
  xtick={0,4.7,10,22,47}, xticklabels={0,{4,7},10,22,47},
  grid=major, grid style={draw=gray!25},
  legend pos=south east, legend style={font=\scriptsize},
  tick label style={font=\scriptsize}, label style={font=\small}]
  \addplot[thick,mark=*,color=verdeplaca] coordinates {(0,7.8) (4.7,70.8) (10,69.7) (22,66.8) (47,64.9)};
  \addlegendentry{TLV9152}
  \addplot[thick,mark=square*,color=cobre,densely dashed] coordinates {(0,3.4) (4.7,72.3) (10,68.8) (22,64.6) (47,61.9)};
  \addlegendentry{TLV9062}
  \draw[densely dotted,color=grisnota] (axis cs:-2,60) -- (axis cs:50,60);
  \node[font=\scriptsize,color=grisnota,anchor=north west] at (axis cs:30,58) {\SI{60}{\degree}};
\end{axis}
\end{tikzpicture}
\caption{Margen de fase del sumador contra CF (datos de la tabla~\ref{tab:c03-cf}). Cualquier CF de algunos picofaradios
resuelve el problema; sin CF el circuito queda al borde de oscilar con cualquiera de los dos op-amps.}
\label{fig:c03-pm}
\end{figure}

\begin{decision}[CF = \SI{10}{\pico\farad} C0G (cambio C16)]
\textbf{Qué se decidió.} Agregar C56 (canal 1) y C57 (canal 2): \SI{10}{\pico\farad}, dieléctrico C0G, \SI{50}{\volt},
0603 (\net{CL10C100JB8NNNC}, LCSC C1634), en paralelo con R7+R8 y con R42+R43.

\textbf{Alternativas.} No poner nada (el circuito original del equipo no tenía CF); \SI{4.7}{\pico\farad};
\SI{22}{\pico\farad} o más.

\textbf{Por qué.}
\begin{itemize}
  \item Sin CF el margen es de \SI{3.4}{\degree} (TLV9062) o \SI{7.8}{\degree} (TLV9152): riesgo real de oscilación.
  \item \SI{4.7}{\pico\farad} ya da \SI{70.8}{\degree}, pero el margen depende de $C_{\mathrm{SUM}}$, que tiene una
        parte estimada (la pista). Con \SI{10}{\pico\farad}, $C_F$ pesa más que la incertidumbre y el margen se queda
        cerca de \SI{70}{\degree} aunque la capacidad real del nodo sea algo distinta.
  \item \SI{10}{\pico\farad} limita la banda de la señal a $1/(2\pi R_f C_F) = 1/(2\pi\cdot\SI{99.8}{\kilo\ohm}\cdot
        \SI{10}{\pico\farad}) \approx \SI{160}{\kilo\hertz}$ (el script da \SI{166}{\kilo\hertz} con su resolución de
        frecuencia). A \SI{20}{\kilo\hertz} la atenuación es $1/\sqrt{1+(20/160)^2}$, unos \SI{0.07}{\decibel}:
        nada en la banda de audio. Y ayuda a filtrar ruido de alta frecuencia antes del ADC.
  \item Con \SI{22}{\pico\farad} o \SI{47}{\pico\farad} la banda bajaría a \SI{76}{\kilo\hertz} o \SI{35}{\kilo\hertz}
        sin ganar margen.
  \item C0G porque su capacidad no cambia con la tensión ni con la temperatura (un X7R de \SI{10}{\pico\farad} sí
        cambiaría, y la banda y el margen dependen directamente de ese valor).
\end{itemize}

\textbf{Fuente.} \archivo{tools/analisis_opamp.py}, \archivo{tools/changes.json} (C16), \archivo{CALCULOS.md} §9.
\end{decision}

\begin{ojo}[es un modelo, no una simulación del fabricante]
El modelo de dos polos es una aproximación ajustada a dos números del datasheet. No se simuló con el modelo SPICE de
TI. La prueba real es en el banco: aplicar un escalón (el puente JP1 o un cambio de código del DAC) y mirar
\net{REST_OUT} en TP2 con el osciloscopio: no debe haber oscilación ni sobrepico \medir. El sitio de C56 junto a U3 y
la longitud de las pistas del nodo \net{SUM} se tratan en el capítulo~\ref{cap:ruteo}.
\end{ojo}

% =====================================================================================================
\section{TLV9152 en lugar del TLV9062}\label{sec:c03-tlv}

\subsection{Por qué se cambió}

El diseño del equipo usaba el TLV9062IDDFR (LCSC C2867884). En la auditoría de piezas para JLCPCB tenía \textbf{2
unidades en stock en JLC y 0 en LCSC}, y la placa necesita 2 por tarjeta (U3 y U9)
(\archivo{fabricacion/auditoria_piezas.md}). Se reemplazó por el TLV9152IDDFR (LCSC C2861011), que tenía stock y viene
en el mismo encapsulado.

\begin{table}[htbp]
\centering\small
\caption{TLV9062 frente a TLV9152. Datos del TLV9152 leídos de SBOS986E; los del TLV9062, de SBOS839N tal como los
usa \archivo{tools/analisis_opamp.py}.}
\label{tab:c03-tlv}
\begin{tabular}{@{}L{0.33\linewidth}L{0.27\linewidth}L{0.32\linewidth}@{}}
\toprule
Característica & TLV9062 (original) & TLV9152 (placa)\\
\midrule
Encapsulado y pines & TSOT-23-8 (DDF) & mismo DDF y mismo orden de pines \doc{SBOS986E fig. 5-4 p.6}\\
Alimentación & — & \SIrange{2.7}{16}{\volt} \doc{SBOS986E p.1}\\
Entrada rail-to-rail & sí & $(V-)-\SI{0.1}{\volt}$ a $(V+)+\SI{0.1}{\volt}$ \doc{SBOS986E p.11}\\
GBW & \SI{10}{\mega\hertz} \doc{SBOS839N p.13} & \SI{4.5}{\mega\hertz} \doc{SBOS986E p.12}\\
Margen de fase a G = +1 & \SI{55}{\degree} & \SI{60}{\degree} ($R_L = \SI{10}{\kilo\ohm}$)\\
Carga capacitiva admitida & — & \SI{1000}{\pico\farad} \doc{SBOS986E p.12}\\
Capacidad de entrada & — & $C_{ID}$ = \SI{3}{\pico\farad}, $C_{ICM}$ = \SI{1}{\pico\farad} \doc{SBOS986E p.11}\\
Consumo por amplificador & — & \SI{560}{\micro\ampere} típ., \SI{685}{\micro\ampere} máx. \doc{SBOS986E p.12}\\
Ruido de tensión & — & \SI{10.8}{\nano\volt\per\sqrt\hertz} (auditoría de piezas)\\
Stock (noviembre 2026) & 2 en JLC, 0 en LCSC & 542 en JLC\\
\bottomrule
\end{tabular}
\end{table}

Los guiones son datos del TLV9062 que no se leyeron para este informe. Lo importante del cambio: el pinout es el mismo
(el footprint y el esquemático no cambian) y el GBW baja a menos de la mitad. Por eso hubo que rehacer el análisis de
estabilidad con el TLV9152 (tabla~\ref{tab:c03-cf}) y además revisar dos cosas nuevas.

\subsection{Comprobación 1: la otra mitad (el seguidor)}

Un seguidor es la configuración más exigente para la estabilidad de un op-amp (ganancia 1, toda la salida vuelve a la
entrada). Lo que lo pone en riesgo es la \emph{carga capacitiva} en la salida.

\begin{calculo}[carga del seguidor U3B]
\begin{itemize}
  \item La red \net{VSHIFT} solo tiene U3.3, U3.6 y U3.7 (\archivo{tools/target.json}): no hay condensadores.
  \item La carga es la entrada $+$ de U3A: \SI{3}{\pico\farad} diferenciales + \SI{1}{\pico\farad} de modo común
        \doc{SBOS986E p.11}, más la pista: unos \SI{5}{\pico\farad}.
  \item El TLV9152 admite \SI{1000}{\pico\farad} de carga y tiene \SI{60}{\degree} de margen a ganancia 1
        \doc{SBOS986E p.12}.
  \item El filtro R12/C4 está en la \emph{entrada} del seguidor (\net{DAC_F}), fuera de su lazo, así que no afecta su
        estabilidad.
\end{itemize}
Resultado: \SI{5}{\pico\farad} frente a \SI{1000}{\pico\farad} admitidos. Estable por datos del datasheet; falta verlo
en el osciloscopio \medir. El canal 2 (U9B con R37/C51) es idéntico.
\end{calculo}

\subsection{Comprobación 2: la región de transición de la entrada}\label{sec:c03-transicion}

\begin{idea}[dos pares de entrada]
Para que un op-amp tenga entrada rail-to-rail, por dentro tiene \emph{dos} pares diferenciales en paralelo: un par de
transistores P que funciona desde un poco debajo de $V-$ hasta cerca de $V+$ menos unos voltios, y un par N que
funciona cerca de $V+$. En la zona donde se pasa de uno al otro trabajan los dos a la vez, y ahí el op-amp es peor: su
offset cambia, y empeoran el rechazo de modo común (CMRR), el de la fuente (PSRR), el ruido y la distorsión.
\end{idea}

En el TLV9152 el par P funciona hasta unos $(V+) - \SI{2}{\volt}$ y el par N desde unos $(V+) - \SI{1}{\volt}$; entre
los dos está la región de transición, cuyos límites son \emph{típicos} y varían con el proceso de fabricación
\doc{SBOS986E §7.3.4 p.26}. La CMRR de la tabla eléctrica solo está garantizada por debajo de $(V+) - \SI{2}{\volt}$
\doc{SBOS986E p.11}.

¿Qué tensión ven las entradas? En el sumador, las dos entradas de U3A están a \net{VSHIFT} (la $+$ directamente y la
$-$ por ser tierra virtual). En el seguidor, las dos están a \net{DAC_F} $\approx$ \net{VSHIFT}. O sea: \textbf{el modo
común de las dos mitades es \net{VSHIFT}}.

\begin{calculo}[dónde empieza la región con nuestro riel]
$V+$ es \net{5V_A}: \SI{5}{\volt} $\pm\SI{1.5}{\percent}$ del buck \doc{SNVSB49D p.8} menos la caída en R33
(\SI{10}{\ohm}) con \SIrange{10.8}{14}{\milli\ampere} de consumo:
\begin{align*}
  V+_{\min} &= 5\cdot 0{,}985 - 0{,}014\cdot 10 = \SI{4.785}{\volt}, &
  V+_{\max} &= 5\cdot 1{,}015 - 0{,}0108\cdot 10 = \SI{4.967}{\volt}.
\end{align*}
La región de transición empieza en $V+ - \SI{2}{\volt}$: entre \SI{2.785}{\volt} y \SI{2.967}{\volt}.
\end{calculo}

Con la ecuación~\eqref{eq:c03-vshift} se puede calcular \net{VSHIFT} para cada objetivo de \net{REST_OUT} y cada
polarización $V_B$ del micrófono. El DAC además tiene un tope: unos \SI{3.21}{\volt} en el peor caso (3V3 al
$-\SI{1.5}{\percent}$, menos la caída en R9, por $4095/4096$). La tabla~\ref{tab:c03-transicion} es la salida del
script.

\begin{table}[htbp]
\centering\small
\caption{\net{VSHIFT} necesario según el objetivo de \net{REST_OUT} y la polarización del micrófono
(\archivo{tools/analisis_opamp.py}, función \texttt{transicion()}). ``Par P'' = fuera de la región de transición.}
\label{tab:c03-transicion}
\begin{tabular}{@{}lllll@{}}
\toprule
Objetivo de \net{REST_OUT} & $V_B$ = \SI{8}{\volt} & \SI{10}{\volt} & \SI{12.5}{\volt} (RG-50) & \SI{14}{\volt}\\
\midrule
\SI{1.50}{\volt} & \SI{2.09}{\volt} par P & \SI{2.27}{\volt} par P & \SI{2.50}{\volt} par P & \SI{2.63}{\volt} par P\\
\SI{2.00}{\volt} & \SI{2.54}{\volt} par P & \SI{2.73}{\volt} par P & \textcolor{cobre}{\SI{2.95}{\volt} transición} & \textcolor{cobre}{\SI{3.09}{\volt} transición}\\
\SI{2.45}{\volt} (mitad del riel) & \textcolor{cobre}{\SI{2.95}{\volt} transición} & \textcolor{cobre}{\SI{3.13}{\volt} transición} & \textcolor{rojoojo}{DAC al tope} & \textcolor{rojoojo}{DAC al tope}\\
\bottomrule
\end{tabular}
\end{table}

Lee la última fila con atención: la idea ``natural'' de centrar la salida a mitad del riel
(\SI{2.45}{\volt}) mete \net{VSHIFT} en la región de transición, y con el RG-50 el DAC ni siquiera alcanza la tensión
necesaria (pediría más de \SI{3.21}{\volt}).

\begin{decision}[objetivo de REST\_OUT $\approx$ \SI{1.5}{\volt} (pendiente de aprobación)]
\textbf{Qué se recomienda.} Que el firmware lleve \net{REST_OUT} a unos \SI{1.5}{\volt} (\net{DC_MON}
$\approx$ \SI{0.75}{\volt}). Es una recomendación: el repositorio no tiene un objetivo fijado y falta la aprobación
del dueño.

\textbf{Alternativas.} Mitad del riel (\SI{2.45}{\volt}); \SI{2}{\volt}; \SI{1.3}{\volt}.

\textbf{Por qué.}
\begin{itemize}
  \item Con \SI{1.5}{\volt}, \net{VSHIFT} queda entre \SI{2.09}{\volt} y \SI{2.63}{\volt} para polarizaciones de
        \SIrange{8}{14}{\volt}: siempre en el par P, fuera de la región.
  \item La salida conserva margen: $1{,}5 \pm 0{,}32$~V con el RG-50 a \SI{130}{\decibel} pico, y $1{,}5 \pm 1{,}08$~V
        con la excursión máxima del SV 12. Lo más bajo es \SI{0.42}{\volt}, lejos de GND.
  \item El margen más chico es con $V_B = \SI{14}{\volt}$: \SI{2.63}{\volt} contra \SI{2.785}{\volt}, solo
        \SI{155}{\milli\volt}. Como el límite de la región es típico, si se quiere más margen se puede bajar a unos
        \SI{1.3}{\volt}; el piso lo pone la excursión de \SI{1.08}{\volt} del SV 12, cuya polarización aún no se midió.
\end{itemize}
\textbf{¿Qué tan grave sería caer en la región?} Aquí el modo común es un nivel de DC que no se mueve con la señal
(configuración inversora y lazo de DC). El cambio de offset es DC y el lazo lo corrige; la peor CMRR casi no importa
porque el modo común no varía. Lo que el lazo no corrige es el aumento de ruido y de distorsión. Por eso conviene
evitarla, aunque no sea un desastre.

\textbf{Fuente.} \archivo{CALCULOS.md} §9, \archivo{tools/analisis_opamp.py}, SBOS986E §7.3.4.
\end{decision}

% =====================================================================================================
\section{El canal 2 es un espejo del canal 1}\label{sec:c03-ch2}

El Altium del equipo solo tenía dibujado el canal 1. El canal 2 se generó copiando pieza por pieza y cambiando el
nombre de las redes (\archivo{tools/design.py}, diccionarios \texttt{MIRROR} y \texttt{NETMAP}). La
tabla~\ref{tab:c03-espejo} da la correspondencia.

\begin{table}[htbp]
\centering\small
\caption{Correspondencia entre el canal 1 y el canal 2.}
\label{tab:c03-espejo}
\begin{tabular}{@{}L{0.30\linewidth}L{0.28\linewidth}L{0.34\linewidth}@{}}
\toprule
Función & Canal 1 & Canal 2\\
\midrule
Fuente IEPE & U1, R1, C1–C3, R2$parallel$R3, R4–R6 & U8, R36, C48–C50, R35, R34\\
Escalón & R47 + JP1 & R48 + JP2\\
Conector y punto de prueba & J2, TP1, habilitación J1 & J8, TP3, habilitación J7\\
Cadena de entrada & R13, R14, R16, R18 & R38, R39, R40, R41\\
Realimentación y CF & R7, R8, C56 & R42, R43, C57\\
Op-amp & U3 (A y B) & U9 (A y B)\\
DAC y filtro & U2 (0x60), R12, C4 & U10 (0x61), R37, C51\\
Alimentación del DAC & R9, C5, C6 & R46, C52, C53\\
\net{DC_MON} & R15, R17, C8 & R44, R45, C55 (\net{DC_MON2}, PA1)\\
Desacoplo del op-amp & C7 & C54\\
Redes & \net{IEPE}, \net{SUM}, \net{VSHIFT}, \net{REST_OUT} & \net{IEPE_SV}, \net{SUM_SV}, \net{VSHIFT_SV}, \net{REST_SV}\\
\bottomrule
\end{tabular}
\end{table}

Las únicas diferencias eléctricas son:
\begin{itemize}
  \item la fuente: \SI{1.51}{\milli\ampere} / \SI{2.51}{\milli\ampere} en vez de \SI{2.81}{\milli\ampere} /
        \SI{3.98}{\milli\ampere}, con una sola $R_{\mathrm{SET}}$ (R34) y una sola $R_{\mathrm{OUT}}$ (R35);
  \item la dirección del DAC: el pin A0 de U10 va a \net{3V3_DAC_SV} (0x61) en vez de GND;
  \item la salida va a la entrada derecha del ADC (\net{REST_SV} $\to$ C14 $\to$ \net{AC_RIGHT}).
\end{itemize}
En la placa, el canal 2 también quedó con la misma geometría que el canal 1 desplazada \SI{22.5}{\milli\meter}
(capítulo~\ref{cap:placement}).

% =====================================================================================================
\section{Alimentación y desacoplo del canal}\label{sec:c03-alim}

\subsection{El riel 5V\_A}

Los op-amps U3 y U9 (y el PCM1808) no se alimentan del \SI{5}{\volt} del buck directamente, sino de \net{5V_A}, que
sale de un filtro RC: R33 (\SI{10}{\ohm}) en serie y C47 (\SI{10}{\micro\farad}) a tierra.
\begin{equation*}
  f_c = \frac{1}{2\pi\cdot\SI{10}{\ohm}\cdot\SI{10}{\micro\farad}} \approx \SI{1.6}{\kilo\hertz}.
\end{equation*}
A \SI{1}{\mega\hertz}, la frecuencia del buck, eso atenúa unas $10^6/1{,}6\cdot10^3 \approx 600$ veces
(\SI{56}{\decibel}) en la aproximación de primer orden. Además los op-amps tienen su propio rechazo de fuente. El costo
es la caída en R33 (\SIrange{0.11}{0.14}{\volt}), que ya se usó arriba para calcular dónde empieza la región de
transición. Los detalles del buck y del filtro están en el capítulo~\ref{cap:alim}.

\subsection{Desacoplos}

\begin{itemize}
  \item \textbf{C7} (\SI{100}{\nano\farad} X7R) en el pin 8 de U3 (\net{5V_A}); \textbf{C54} igual en U9. Un op-amp
        necesita su condensador lo más cerca posible del pin de alimentación: es la ``reserva'' de carga para los
        picos de corriente rápidos que la pista no alcanza a entregar. En la placa C7 quedó a \SI{0.975}{\milli\meter}
        de U3.8 y C54 a \SI{0.74}{\milli\meter} de U9.8 \medido.
  \item \textbf{C5} (\SI{10}{\micro\farad}) y \textbf{C6} (\SI{100}{\nano\farad}) en el VDD del DAC (\net{3V3_DAC});
        C52/C53 en el canal 2. El datasheet del MCP4725 pide \SI{0.1}{\micro\farad} (más \SI{10}{\micro\farad}) a no
        más de \SI{4}{\milli\meter} de VDD \doc{DS22039D §8.3}; quedaron entre \SI{1.22}{\milli\meter} y
        \SI{1.48}{\milli\meter} \medido.
  \item Todos los \SI{100}{\nano\farad} del canal son X7R (cambio C7): los Y5V del diseño original pierden buena parte
        de su capacidad con la tensión y la temperatura.
\end{itemize}

Cómo se colocaron y rutearon esos condensadores (con su vía a GND pegada al pad) se explica en el
capítulo~\ref{cap:ruteo}, y las excepciones aceptadas en la distancia de algunos de ellos en el
capítulo~\ref{cap:banco}.

% =====================================================================================================
\section{Qué queda por comprobar en el banco}\label{sec:c03-banco}

\begin{itemize}
  \item Corriente de cada fuente con amperímetro en J1/J7, con y sin puente: \
um{2.81}/\SI{3.98}{\milli\ampere} y
        \
um{1.51}/\SI{2.51}{\milli\ampere} \medir.
  \item Polarización del SV 12 a \SI{2.5}{\milli\ampere} (esperado entre \
um{7} y \SI{14}{\volt}) y si le alcanza el
        compliance para su excursión máxima \medir.
  \item \net{VSHIFT} y \net{REST_OUT} en DC con cada micrófono; que \net{VSHIFT} quede por debajo de unos \SI{2.7}{\volt}
        con el objetivo de \SI{1.5}{\volt} \medir.
  \item Respuesta a un escalón (puente JP1/JP2 o código del DAC) en TP2/TP4: sin oscilación ni sobrepico \medir.
  \item Ruido de la cadena completa con el micrófono en una caja silenciosa \medir.
\end{itemize}

% =====================================================================================================
\section*{Resumen}
\addcontentsline{toc}{section}{Resumen}
\begin{itemize}
  \item El LT3092 hace una fuente de corriente con $I = \SI{10}{\micro\ampere}(1 + R_{\mathrm{SET}}/R_{\mathrm{OUT}})$:
        \SI{2.81}{\milli\ampere} en el canal 1 (RG-50) y \SI{1.51}{\milli\ampere} en el canal 2 (SV 12). Los puentes
        JP1/JP2 dan un escalón a \
um{3.98} y \SI{2.51}{\milli\ampere} para la demostración del lazo.
  \item Con \SI{24}{\volt} quedan unos \SI{6}{\volt} de margen para el RG-50; el SV 12 hay que medirlo.
  \item El sumador U3A da $V_{\mathrm{REST}} = 1{,}0997\,V_{\mathrm{SHIFT}} - 0{,}0997\,v_{\mathrm{IEPE}}$: resta la
        polarización en DC y atenúa la señal diez veces, para que la excursión del SV 12 (\SI{10.8}{\volt}) quepa en los
        $\pm\SI{1.5}{\volt}$ del ADC.
  \item \net{VSHIFT} lo genera un DAC de 12 bits (0x60/0x61) con su alimentación filtrada, un RC de \SI{1}{\milli\second}
        y un seguidor. El firmware mide \net{DC_MON} y corrige \net{VSHIFT} con $e/1{,}1$; el lazo queda quieto
        mientras se graba.
  \item El nodo \net{SUM} (\SI{1}{\mega\ohm} con \SI{5.5}{\pico\farad}) dejaba el sumador con menos de
        \SI{8}{\degree} de margen. CF = \SI{10}{\pico\farad} C0G lo sube a unos \SI{70}{\degree} y limita la banda a
        unos \SI{160}{\kilo\hertz}.
  \item El TLV9152 reemplazó al TLV9062 por stock: mismo pinout, menos GBW. El seguidor es estable (\
um{5} de
        \SI{1000}{\pico\farad}). Para no caer en la región de transición de la entrada se recomienda un objetivo de
        \net{REST_OUT} de unos \SI{1.5}{\volt}, pendiente de aprobación.
  \item El canal 2 es un espejo exacto del canal 1, con otra corriente y otra dirección I\textsuperscript{2}C.
\end{itemize}
